\documentclass{article}
\usepackage{amsmath}
\usepackage{amssymb}

\begin{document}

\[ \text{Exercice 2} \]
\[ \text{a) } \sqrt{x+6} = 2 - \sqrt{x+1} \quad (...)^2 \]
\[ \implies x+6 = (2 - \sqrt{x+1})^2 \quad D(f) = ]-\infty,-6]U[-1,\infty]\]
\[ x+6 = 4 - 2\sqrt{x+1} + (x+1) \]
\[ \implies x+6 = x+5 - 2\sqrt{x+1} \]
\[ \implies 6 = 5 -2\sqrt{x+1} \]
\[ \implies 1 = -2\sqrt{x+1} \]
\[ \implies -\frac{1}{2} = \sqrt{x+1} \quad (...)^2\]
\[ \frac{1}{4} = x+1 \]
\[ \implies x = \frac{1}{4} - \frac{1}{1} = \frac{-3}{4} \]
\[ x \ge -6 \text{ dans domaine de definition} \]
\[ x \ge -1 \]
\[ \text{b) } \log_2(x-2) + \log_2(10-x) = 3 + \log_2(8-x) \]
\[ \log_2((x-2)(10-x)) = 3 + \log_2(8-x) \]
\[ \log_2\left(\frac{(x-2)(10-x)}{8-x}\right) = 3 \]
\[ (x-2)(10-x) = 10x - x^2 - 20 + 2x \implies -x^2 + 12x - 20 \]
\[ x^2+-2x+20 = (x-10)(x-2) \]
\[ (x-4)+12 (reste) \]
\[ \log_2((x-4)+12) = 3 \]

\end{document}